\chapter{Calibrazione, Triangolazione e Geometria Proiettiva}

\section{Introduzione e Motivazione}
Nelle lezioni precedenti è stato introdotto il modello geometrico non lineare della telecamera. L'obiettivo principale di questa trattazione è formalizzare rigorosamente tale modello e affrontare il problema fondamentale della \textbf{calibrazione della telecamera}.

Calibrare una telecamera significa stimare i parametri che ne governano il modello geometrico (distanza focale, centro principale, parametri intrinseci ed estrinseci). La conoscenza approfondita di tali parametri consente di formalizzare compiutamente le leggi di proiezione geometrica, rendendo il modello utilizzabile per compiti complessi come la ricostruzione 3D, la robotica e la visione artificiale.

\section{Modello di Proiezione, Raggi Ottici e Triangolazione}

\subsection{Proiezione e Ambiguità della Singola Vista}
In un modello stilizzato bidimensionale, la telecamera può essere rappresentata da:
\begin{itemize}
    \item Un \textbf{piano immagine} (rappresentato come un segmento nel piano);
    \item Un \textbf{centro di proiezione} (o centro ottico $O$).
\end{itemize}

La proiezione di un punto $P$ dello spazio sul piano immagine avviene lungo il \textbf{raggio ottico} passante per $P$ e per il centro $O$. Tuttavia, una singola immagine presenta un’ambiguità intrinseca di profondità: qualsiasi punto appartenente al medesimo raggio ottico produce la stessa identica proiezione sul piano immagine.

Calibrare la telecamera significa essere in grado di ricostruire la retta orientata dello spazio che costituisce il raggio ottico a partire dalle coordinate discrete del punto immagine. Per fare ciò è necessario conoscere la posizione del centro ottico e parametri fondamentali come la distanza focale.

\subsection{Triangolazione a Due o Più Viste}
Per superare l'ambiguità di profondità e ricostruire le coordinate tridimensionali del punto $P$, si impiega un secondo punto di vista:
\begin{enumerate}
    \item Si dispone di una seconda telecamera calibrata;
    \item Si individuano due punti immagine \textbf{corrispondenti} (ovvero generati dallo stesso punto fisico $P$ nello spazio);
    \item Si tracciano i due raggi ottici associati ai punti corrispondenti dai rispettivi centri di proiezione.
\end{enumerate}

L'intersezione tridimensionale dei due raggi ottici individua univocamente la posizione di $P$. Questo processo prende il nome di \textbf{triangolazione}, poiché individua un triangolo in cui la base è il segmento che unisce i due centri ottici e i lati sono i due raggi di proiezione.

\subsection{Gestione dell'Errore e Stima ai Minimi Quadrati}
Nel mondo reale, i due raggi ottici non si intersecano mai perfettamente a causa di molteplici sorgenti di errore:
\begin{itemize}
    \item Errore nella stima dei parametri di calibrazione delle telecamere;
    \item Errore indotto dalla discretizzazione delle coordinate sul sensore (lavoro in pixel);
    \item Rumore di acquisizione delle immagini.
\end{itemize}

Per mitigare l'impatto degli errori, si calcola il punto di minima distanza tra le due rette sgembe (ad esempio il punto medio del segmento di minima distanza). Aumentando il numero di immagini ($3, 4, 5$ o più punti di vista), è possibile formulare uno \textbf{stimatore ai minimi quadrati} (Least Squares) che media gli errori commessi, fornendo una stima ottimale della posizione 3D del punto.

\section{Autocalibrazione e Struttura dal Movimento (SfM)}
Un'applicazione di grande rilevanza pratica (ad esempio in ambito architettonico o nel rilievo territoriale tramite droni) riguarda il caso in cui non si dispone di telecamere preventivamente calibrate. 

In questi contesti si utilizzano algoritmi di \textbf{autocalibrazione} e tecniche di \textit{Structure from Motion} (SfM). Acquisendo un'ampia sequenza di immagini da punti di vista differenti, si sfruttano vincoli geometrici ridondanti per ottimizzare contemporaneamente sia le posizioni dei punti 3D della scena sia i parametri intrinseci ed estrinseci delle telecamere.

\section{Limiti della Geometria Affine e Necessità dello Spazio Proiettivo}

\subsection{Invarianti Affini e Decomposizione SVD}
Nello spazio euclideo o affine $\mathbb{R}^n$, una generica trasformazione affine ha la forma:
\begin{equation}
    y = A x + b
\end{equation}
con $A$ matrice non singolare (invertibile) e $b$ vettore di traslazione.

Sfruttando la decomposizione ai valori singolari (\textit{Singular Value Decomposition}, SVD), la matrice $A$ può essere espressa come:
\begin{equation}
    A = U D V^T
\end{equation}
dove $U$ e $V$ sono matrici ortogonali (che rappresentano rotazioni o riflessioni) e $D$ è una matrice diagonale contenente i valori singolari ($\sigma_1, \sigma_2, \dots$). 

L'azione della matrice $D$ consiste nello scalare lo spazio lungo gli assi con fattori $\sigma_i$ potenzialmente differenti. L'effetto geometrico complessivo di una trasformazione affine include quindi rotazioni, traslazioni e deformazioni (dilatazioni anisotropiche).

Tuttavia, le trasformazioni affini preservano una fondamentale proprietà invariante: \textbf{il parallelismo}. Rette parallele nello spazio di partenza rimangono strettamente parallele anche dopo una trasformazione affine.

\subsection{I Punti di Fuga e la Proiezione Prospettica}
La proiezione prospettica effettuata da una telecamera reale \textbf{non preserva il parallelismo}: rette parallele nel mondo 3D proiettano sul piano immagine in rette incidenti che convergono in un punto detto \textbf{punto di fuga}.

Geometricamente, il punto di fuga non rappresenta la proiezione di un punto al finito dello spazio, bensì la proiezione della \textbf{direzione} della retta. Poiché la geometria affine non consente di rappresentare tali trasformazioni né le direzioni come elementi dello spazio punti, si rende necessario introdurre una nuova struttura geometrica: lo \textbf{spazio proiettivo}.

\section{Formalizzazione dello Spazio Proiettivo e Coordinate Omogenee}

\subsection{Costruzione dello Spazio Proiettivo $\mathbb{P}^n$}
Per rappresentare sia i punti al finito sia le direzioni (punti all'infinito), estendiamo lo spazio vettoriale $\mathbb{R}^{n+1}$ rimuovendo il vettore nullo: $\mathbb{R}^{n+1} \setminus \{0\}$. Si rimuove l'origine poiché essa corrisponde al centro di proiezione, punto non proiettabile sul piano immagine.

Definiamo su $\mathbb{R}^{n+1} \setminus \{0\}$ la seguente relazione di equivalenza $\sim$:
\begin{equation}
    X \sim Y \iff \exists \lambda \in \mathbb{R} \setminus \{0\} \quad \text{tale che} \quad Y = \lambda X
\end{equation}

Ogni classe di equivalenza identifica una retta vettoriale passante per l'origine (privata dell'origine stessa). Lo \textbf{spazio proiettivo $n$-dimensionale} $\mathbb{P}^n$ è definito come lo spazio quoziente:
\begin{equation}
    \mathbb{P}^n = \frac{\mathbb{R}^{n+1} \setminus \{0\}}{\sim}
\end{equation}

\subsection{Coordinate Omogenee}
Un punto $P \in \mathbb{P}^n$ è individuato da una tupla di $n+1$ elementi $[x_1, x_2, \dots, x_n, x_{n+1}]^T \neq 0$. Tale rappresentazione è definita \textbf{a meno di un fattore di scala non nullo} $\lambda$:
\begin{equation}
    [x_1, \dots, x_{n+1}]^T \equiv [\lambda x_1, \dots, \lambda x_{n+1}]^T, \quad \forall \lambda \neq 0
\end{equation}

Tali coordinate prendono il nome di \textbf{coordinate omogenee}.

\section{Base di uno Spazio Proiettivo e il Punto Unità}

\subsection{Il Problema delle Combinazioni Lineari nello Spazio Proiettivo}
In uno spazio vettoriale $\mathbb{R}^{n+1}$, una base è costituita da $n+1$ vettori linearmente indipendenti. Nello spazio proiettivo, tuttavia, i punti sono definiti a meno di scala. 

Se provassimo a generare un punto proiettivo mediante combinazione lineare dei rappresentanti di una base vettoriale $B_1, \dots, B_{n+1}$, cambiando arbitrariamente i fattori di scala dei singoli rappresentanti ($B_i \to \gamma_i B_i$) otterremmo punti proiettivi tra loro non equivalenti, facendo crollare l'unicità della rappresentazione rispetto alla base.

\subsection{Definizione di Riferimento Proiettivo e Punto Unità}
Per risolvere questa ambiguità e fissare univocamente i fattori di scala relativi tra gli elementi della base, è necessario introdurre un ulteriore punto, detto \textbf{punto unità} $B_{n+2}$.

Un \textbf{riferimento (o base) proiettivo} per $\mathbb{P}^n$ è costituito da un insieme di $n+2$ punti $\{B_1, B_2, \dots, B_{n+1}, B_{n+2}\}$ tali che:
\begin{itemize}
    \item I primi $n+1$ punti sono proiettivamente indipendenti (i loro rappresentanti vettoriali in $\mathbb{R}^{n+1}$ sono linearmente indipendenti);
    \item Il punto $B_{n+2}$ (punto unità) vincola la scelta dei rappresentanti vettoriali $b_i$ di $B_i$ imponendo la condizione:
    \begin{equation}
        b_{n+2} = b_1 + b_2 + \dots + b_{n+1}
    \end{equation}
\end{itemize}

Vincolando la somma dei rappresentanti al punto unità (ad esempio il punto di coordinate omogenee $[1, 1, \dots, 1]^T$), la scelta dei rappresentanti della base è determinata a meno di un unico fattore di scala globale $\lambda$ comune a tutti gli elementi, preservando la consistenza della struttura proiettiva.

\section{Immersione dello Spazio Affine e Interpretazione Geometrica}

\subsection{Punti al Finito e Punti all'Infinito}
Sia $\mathbb{R}^n$ lo spazio affine di partenza. Possiamo immergere $\mathbb{R}^n$ in $\mathbb{P}^n$ tramite la mappa canonica:
\begin{equation}
    (x_1, x_2, \dots, x_n) \longmapsto [x_1, x_2, \dots, x_n, 1]^T
\end{equation}

Considerando un generico punto di $\mathbb{P}^n$ di coordinate $[x_1, \dots, x_n, x_{n+1}]^T$:
\begin{itemize}
    \item Se $x_{n+1} \neq 0$, il punto appartiene all'immagine dello spazio affine (\textbf{punto al finito}). Le sue coordinate affini si ottengono de-omogeneizzando:
    \begin{equation}
        \left( \frac{x_1}{x_{n+1}}, \frac{x_2}{x_{n+1}}, \dots, \frac{x_n}{x_{n+1}} \right)
    \end{equation}
    \item Se $x_{n+1} = 0$, il punto ha coordinata omogenea finale nulla: $[x_1, \dots, x_n, 0]^T$. Tale punto non corrisponde ad alcun punto reale dello spazio affine $\mathbb{R}^n$ ed è definito \textbf{punto all'infinito} (o \textbf{punto ideale}).
\end{itemize}

\subsection{Interpretazione Geometrica in $\mathbb{P}^2$}
Consideriamo la costruzione del piano proiettivo $\mathbb{P}^2$ a partire dallo spazio tridimensionale $\mathbb{R}^3$:
\begin{enumerate}
    \item I punti di $\mathbb{P}^2$ corrispondono alle rette passanti per l'origine in $\mathbb{R}^3$.
    \item Consideriamo un piano affine $\Pi$ non passante per l'origine, ad esempio il piano $z = 1$.
    \item Tutte le rette per l'origine con terza coordinata non nulla ($z \neq 0$) intersecano il piano $\Pi$ in un unico punto di coordinate $(x/z, y/z, 1)$, rappresentando i punti al finito di $\Pi$.
    \item Le rette per l'origine appartenenti al piano $z = 0$ (parallele al piano $\Pi$) non intersecano mai $\Pi$ al finito. Ognuna di queste rette individua una direzione di fascio di rette parallele su $\Pi$: esse costituiscono i **punti all'infinito** di $\mathbb{P}^2$.
\end{enumerate}

L'insieme di tutti i punti all'infinito con $z = 0$ forma la \textbf{retta all'infinito} (o retta ideale) del piano proiettivo. In questo modo, lo spazio proiettivo unifica elegantemente la trattazione dei punti finiti e delle direzioni dello spazio.